\chapter{الأعمال المشابهة}

\begin{chaptersummary}
يستعرض هذا الفصل المنصات والأبحاث القائمة في مجال التعليم الإلكتروني ومؤتمرات
الفيديو التعليمية، وينتهي بتحديد الفجوة التي يسعى المشروع إلى سدّها.
\end{chaptersummary}

\section{مقدمة}
[مقدمة الفصل.]

\section{التعليم والتعلم عن بعد \en{(E-Learning and Distance Education)}}

\subsection{مفهوم التعليم مقابل التعلم \en{(Teaching vs. Learning)}}
[المحتوى من المسودة.]

\subsection{مزايا وعيوب التعليم الإلكتروني}
[المحتوى من المسودة.]

\section{مؤتمرات الفيديو في التعليم}

\subsection{التحديات التي تواجه مؤتمرات الفيديو التعليمية}
[المحتوى من المسودة.]

\section{منصات الاتصال التجارية}

\subsection{أدوات الأغراض العامة \en{(Zoom, Google Meet, Microsoft Teams)}}
[المحتوى من المسودة.]

\subsection{منصات التعليم الإلكتروني المتخصصة \en{(BigBlueButton)}}
[المحتوى من المسودة.]

\section{البحث الأكاديمي في الفصول الدراسية المدعومة بالذكاء الاصطناعي}

\subsection{التوليد الآلي للمحتوى في التعليم}
[المحتوى من المسودة.]

\subsection{المساعد التعليمي المعتمد على \en{RAG}}
[المحتوى من المسودة.]

\section{الفجوة}
% NOTE: numbered "5.4" in the Word draft, which places it under chapter 4.
% It belongs to chapter 3 — LaTeX numbers it correctly here.
[تحديد الفجوة البحثية والتطبيقية.]

\section{الخاتمة}
[خاتمة الفصل.]

\transition{بعد تحديد الفجوة التي يعالجها المشروع، ننتقل في الفصل التالي إلى
تحليل متطلبات النظام وحالات استخدامه.}
